\begin{pastexam}[label={px:ntu-114-dsa-18}]{臺大資網所 114 學年度 第 18 題}
（問答）

(5 points) Consider the task of incrementing a $k$-bit binary counter from all zeroes to all ones. Let $A$ be an
array of length $k$ that simulates the counter, where $A[0]$ represents the least significant bit and $A[k-1]$
represents the most significant bit. We implement the \textsc{Increment}() operation as follows:
\begin{center}
\begin{tabular}{l}
\hline
\textbf{function} \textsc{Increment}() \\
\quad $i=0$ \\
\quad \textbf{while} $A[i]==1$ \textbf{do} \\
\quad\quad $A[i]=0$ \\
\quad\quad $i=i+1$ \\
\quad $A[i]=1$ \\
\hline
\end{tabular}
\end{center}
Unfortunately, this counter is implemented on a deficient random access machine (RAM). On this machine, each
read-access to a single array entry takes 1 unit of time. However, a write to an entry may take different amount
of time. In particular, for any index $i$ ($0\le i<k$), if $i$ is a multiple of 3 but not a multiple of 5, the
memory stick will fizz and it will take exactly $3^{i/3}$ units of time to write the data into $A[i]$. If $i$ is a
multiple of 5 but not a multiple of 3, the memory stick will buzz and it will take exactly $5^{i/5}$ units of time
to write the data into $A[i]$. Writing into all other entries not mentioned above takes 1 unit of time as usual.

Please design a potential function $\Phi(A)$ such that the amortized time (using $\Phi(A)$) per
\textsc{Increment}() operation is still $O(1)$ when incrementing from zero. Make sure your potential function is
rigurously defined, can be evaluated into an actual number by human beings, and as clean and clear without
ambiguity. For example, if you left any undetermined variable (such as ``for some constant $c$'') in your
expression, you \textbf{must} state an actual number for that variable (e.g., ``$c=1$''). You do \textbf{not} need
to verify your answer in the answer sheet. This problem has no partial score.
\end{pastexam}

\begin{pastsoln}
寫 $w(i)$ 表示寫入 $A[i]$ 的時間：$i$ 是 $3$ 的倍數但不是 $5$ 的倍數時
$w(i) = 3^{i/3}$；$i$ 是 $5$ 的倍數但不是 $3$ 的倍數時 $w(i) = 5^{i/5}$；其他
$w(i) = 1$（包括 $i = 0$、$15$、$30$ 這些同時是兩者倍數的）。 令
$L_i(A) = \sum_{j=0}^{i-1} A[j]\, 2^j$，也就是計數器最低 $i$ 位表示的數。

\textbf{答案}：
\[
\Phi(A) = \sum_{i=1}^{k-1} \frac{w(i) + 1}{2^{\,i-1}} \, L_i(A) .
\]
式子裡沒有未定的常數。 例如 $k = 4$、$A[3..0] = 0110$：$L_1 = 0$、$L_2 = 2$、
$L_3 = 6$，$w(1) = w(2) = 1$、$w(3) = 3$，$\Phi = 2 \cdot 0 + \frac{2}{2} \cdot 2 +
\frac{4}{4} \cdot 6 = 8$。

\textbf{想法}。 資結書第~3.3~節的位能「$1$ 的個數」在這裡不夠用。 只看
哪些位是 $1$ 的位能（例如每個 $1$ 存 $w(i) + 1$），在把 $A[t]$ 從 $0$ 寫成 $1$ 的那一
次，要付這次寫入的 $w(t)$，還要替它存下以後清回 $0$ 的錢，平攤成本大約是 $2w(t)$，
$t$ 大的時候沒有上界。 所以寫入 $A[t]$ 的錢（不論寫 $0$ 或寫
$1$）必須在\textbf{之前}就存好。 $A[i]$ 每 $2^i$ 次加一才被寫一次（最低 $i$ 位從全
$1$ 進位的時候），$L_i$ 正好記錄「上次寫 $A[i]$ 之後又加了幾次」：每次加一 $L_i$ 加
$1$，進位到第 $i$ 位時歸零。 每次加一替第 $i$ 位存 $(w(i) + 1)/2^{i-1}$，存 $2^i - 1$
次之後就夠付一次讀加寫。

\textbf{驗證}（題目說答案卷上不用寫，這裡寫完整）。 一開始全是 $0$，$\Phi = 0$；每一
項都不是負的，$\Phi \ge 0$，所以資結書定理~3.2 的條件成立。 從全 $0$ 加到全
$1$ 的過程中不會溢位。 設某次加一之前，最低位起有連續 $t$ 個 $1$（$0 \le t \le k -
1$），$A[t] = 0$。
\begin{itemize}
\item 實際成本：\textbf{while} 讀了 $A[0], \ldots, A[t]$ 共 $t + 1$ 次；寫了
$A[0..t-1]$ 為 $0$、$A[t]$ 為 $1$。 $c = \sum_{i=0}^{t} \bigl(1 + w(i)\bigr)$。
\item $1 \le i \le t$：最低 $i$ 位原本全是 $1$，加一後全是 $0$，$L_i$ 從 $2^i - 1$ 變
成 $0$。
\item $i \ge t + 1$：最低 $i$ 位裡沒有進位出去（$A[t]$ 原本是 $0$），$L_i$ 剛好加 $1$。
\end{itemize}
所以平攤成本是
\[
\hat{c} = \bigl(1 + w(0)\bigr)
+ \sum_{i=1}^{t} \bigl(w(i) + 1\bigr)\left(1 - \frac{2^i - 1}{2^{\,i-1}}\right)
+ \sum_{i=t+1}^{k-1} \frac{w(i) + 1}{2^{\,i-1}} .
\]
$w(0) = 1$。 $i \ge 1$ 時 $\frac{2^i - 1}{2^{i-1}} = 2 - 2^{1-i} \ge 1$，中間的每一項
都 $\le 0$。 所以
\[
\hat{c} \le 2 + \sum_{i=1}^{\infty} \frac{w(i) + 1}{2^{\,i-1}} .
\]
這個無窮級數收斂：$5^{1/5} \approx 1.380 < 3^{1/3} \approx 1.442$，所以每個 $i$ 都有
$w(i) \le (3^{1/3})^i$。 令 $r = 3^{1/3}/2 < 0.73$，
\[
\sum_{i \ge 1} \frac{w(i)}{2^{i-1}} \le 2 \sum_{i \ge 1} r^i = \frac{2r}{1 - r}
< \frac{2 \times 0.73}{0.27} < 5.5, \qquad \sum_{i \ge 1} \frac{1}{2^{i-1}} = 2 ,
\]
因此每次 \textsc{Increment} 的平攤成本 $\hat{c} < 2 + 2 + 5.5 = 9.5$，和 $k$ 無關，
是 $O(1)$。 由資結書定理~3.2，從全 $0$ 開始做 $N$ 次加一的實際總時間不超過
$9.5N$。 實際把級數加起來大約是 $5.20$，所以 $\hat{c} \le 7.2$。
\end{pastsoln}
